% !TEX root = chapter-standalone.tex

\section{Sub-things}
\label{sec:subthings-algebra}

\todotext{J: @J: add material}




\todotext{Not clear yet if we need/want to discuss 'subthings' at all. (One consideration: will be want to talk about subcategories?). If we do talk about sub-things here, then maybe we should also briefly discuss sub-preorders in the section on 'new from old' for them?}

\todotext{In 'subthings' examples, also have negative examples that show how things can fail.}

\paragraph{Subsemigroups}

\todotext{\bernina: Add some text and some examples to this subsection (subsemigroup)}

\begin{definition}[Subsemigroup]\label{def:subsemigroup}
    Let~$\sgrpA = \tup{\sgrpAset, \mtimes}$ be a \SY{semigroup}.
    A \maindef{subsemigroup} of~$\sgrpA$ is:

    \constit

    \begin{enumerate}
        \item A subset~$\sgrpBset \setsubseteq \sgrpAset$.
    \end{enumerate}

    \condit

    \begin{enumerate}
        \item The set~$\sgrpBset$ is closed under the composition operation~$\mtimes$ from~$\sgrpA$:
              \begin{equation}\label{eq:subsemigroup-closure}
                  \prfperiod{
                      \sgrpela \setin \sgrpBset \quad \sgrpelb \setin \sgrpBset
                  }{
                      (\sgrpela \mtimes \sgrpelb) \setin \sgrpBset
                  }
              \end{equation}
    \end{enumerate}
\end{definition}

\begin{example}
    \label{exa:evennum-subsemigroup}
    Consider the \SY{semigroup}~$\tup{\natnumbers,+}$ introduced in \cref{exa:natnum-semigroup}.
    We can take the subset of even natural numbers~$\sgrpBset \setsubset \natnumbers$.
    One can show that the sum of two even numbers is always even, satisfying the closure of~$\sgrpBset$ under the composition operation~$+$.
\end{example}



\paragraph{Submonoids}

\todotext{\bernina: Add some text and some examples to this subsection}

\begin{definition}[Submonoids]\label{def:submonoids}
    Let~$\monoidAdefinition$ be a \SY{monoid}.
    A \maindef{submonoid} of~$\monoidA$ is:

    \constit

    \begin{enumerate}
        \item A subset~$\monoidBset \setsubseteq \monoidAset$.
    \end{enumerate}

    \condit

    \begin{enumerate}
        \item The set~$\monoidBset$ is closed under the composition operation~$\mtimes$ from~$\monoidA$:
              \begin{equation}
                  \prfsemi{\monela \setin \monoidBset \quad \monelb \setin \monoidBset}{\monela \mtimes \monelb \setin \monoidBset}
              \end{equation}

        \item The \SY{neutral element}~$\idmon \setin \monoidAset$ is an element of~$\monoidBset$.
    \end{enumerate}
\end{definition}

\begin{example}
    $\tup{\natnumbers, +, 0}$,~$\tup{\wnumbers, +, 0}$, and~$\tup{\ratnumbers, +, 0}$ are all submonoids of~$\tup{\reals, +, 0}$.
\end{example}

\begin{example}
    $\tup{\natnumbers, \cdot, 1}$  is a \SY{submonoid} of~$\tup{\wnumbers, \cdot, 1}$.
\end{example}

\begin{example}
    $\makeset{ \ela \setin \natnumbers \mid \ela \text{ is even } }$ is a \SY{submonoid} of~$\tup{\natnumbers, +, 0}$.
\end{example}



\paragraph{Subgroups}

\todotextjira{454}{\bernina: Add more to subgroups subsection}

\begin{definition}[Subgroup]\label{def:subgroup}
    Let~$\grpA = \tup{\grpAset, \gtimes, \idgrp, \ginv}$ be a group.
    A \maindef{subgroup} of~$\grpA$ is:

    \constit

    \begin{enumerate}
        \item A subset~$\grpBset \setsubseteq \grpAset$.
    \end{enumerate}

    \condit

    \begin{enumerate}
        \item The set~$\grpBset$ is closed under the composition operation~$\gtimes$ from~$\grpA$:
              \begin{equation}
                  \prfsemi{
                      \grpela \setin \grpBset \quad \grpelb \setin \grpBset
                  }{
                      \grpela \gtimes \grpelb \setin \grpBset
                  }
              \end{equation}

        \item The \SY{neutral element}~$\idgrp \setin \grpAset$ is an element of~$\grpBset$;

        \item The set~$\grpBset$ is closed under the inverse operation~$\ginv$ from~$\grpA$:
              \begin{equation}
                  \prfperiod{
                      \grpela \setin \grpBset
                  }{
                      \ginv(\grpela) \setin \grpBset
                  }
              \end{equation}
    \end{enumerate}
\end{definition}





